\chapter{结论}\label{chap:conclusion}

\section{与现有工作的比较}\label{sec:comparison_related}

本文的工作与现有PINN期权定价研究的主要区别体现在以下几个维度：

\textbf{统一性}。现有工作（Dhiman等\cite{dhiman2023pinn}、Bae等\cite{bae2024localvol}、Hainaut等\cite{hainaut2024heston}）均针对单一模型，本文是首个将BSM、CEV、Heston三大模型统一在单一网络中的工作。统一框架的优势不仅在于工程上的简洁（一套训练流程、一个模型文件），更在于三种模型的联合训练带来了互相增益：BSM/CEV的简单结构为Heston提供了良好的初始化，有助于克服Heston二维PDE的收敛困难。

\textbf{交互方式}。现有PINN工具均为纯数值接口，用户需手动指定模型类型和参数。本文引入LLM路由层，支持自然语言输入，是PINN与LLM结合用于期权定价的首次尝试。与PDE-AGENT\cite{liu2025pdeagent}相比，本文的LLM仅作为路由层（参数提取+模型选择），不参与求解过程，因此延迟更低、可靠性更高。

\textbf{约束处理}。本文采用加法参数化输出，以Black-Scholes解析解为基线，这与Zamxaka\cite{zamxaka2023cev}的参数化PINN思路相近，但本文将其推广至多模型统一框架，并结合软掩码实现了模型间的连续插值。

\section{工作总结}\label{sec:summary}

本文提出并实现了一个基于物理信息神经网络（PINN）与大语言模型（LLM）的统一期权定价框架，针对现有工作中模型孤立和交互门槛高两大问题，提出了系统性的解决方案。

在PINN部分，本文的核心贡献是统一PDE算子的设计。通过软掩码机制$\text{mask}=\tanh(\xi/0.05)^2$，将BSM、CEV、Heston三大模型的扩散系数统一为单一参数化形式，由一个9维输入的深层全连接网络同时覆盖三种模型的解映射。加法参数化输出以Black-Scholes解析解为基线，网络学习绝对修正量，解决了乘法参数化在深度虚值区域的退化问题，同时提供了良好的初始化。ICPINN硬约束技术通过Black-Scholes公式的极限行为自然满足终值条件，无需在损失函数中加入额外约束项。

在LLM部分，本文引入DeepSeek API作为智能路由层，通过结构化JSON输出实现可靠的参数提取和模型选择，支持中英文自然语言输入，使非专业用户无需了解模型细节即可完成期权定价。

实验结果验证了统一框架的有效性：
\begin{itemize}
  \item 统一PINN在BSM上达到MAE=0.067（ATM相对误差0.31\%），在Heston上达到MAE=0.031（ATM相对误差0.52\%）；
  \item 推断速度（GPU约2ms）显著优于有限差分法（约500ms）和蒙特卡洛方法（约2000ms）；
  \item LLM路由层在20个测试用例上实现了100\%的模型选择准确率；
  \item 消融实验验证了软掩码、加法参数化和数据锚点三个核心组件的必要性。
\end{itemize}

\section{局限性}\label{sec:limitations}

本文的工作存在以下局限性，需要在未来工作中进一步改进：

\textbf{CEV精度有待提升。}CEV模型的ATM相对误差偏高（37.95\%），部分原因是参考解（Crank-Nicolson有限差分）本身的数值误差，但PINN对CEV的求解质量也有改进空间。未来可以通过增加CEV专用的数据锚点、调整损失权重或引入自适应配点采样来提升CEV精度。

\textbf{参数空间覆盖有限。}当前统一PINN针对固定的参数组合（BSM: $\sigma=0.2$；CEV: $\sigma=0.2,\beta=0.5$；Heston: $\kappa=2,\theta=0.04,\xi=0.3,\rho=-0.7$）训练，对其他参数组合的泛化能力未经验证。真正的参数化PINN需要在参数空间上进行联合训练，这会显著增加训练难度和计算代价。

\textbf{仅支持欧式看涨期权。}当前框架仅针对欧式看涨期权，不支持美式期权（需要处理自由边界问题）、看跌期权（可通过put-call parity转换）或奇异期权（障碍期权、亚式期权等）。

\textbf{LLM依赖外部API。}当前LLM路由层依赖DeepSeek API，需要网络连接，且存在API调用延迟（约500ms至2000ms）。未来可以考虑使用本地部署的小型LLM（如Qwen-7B）替代，降低延迟和依赖。

\textbf{缺乏不确定性量化。}当前框架仅输出点估计，不提供预测不确定性。在金融应用中，不确定性量化（如置信区间）对风险管理至关重要。未来可以引入贝叶斯PINN或集成方法提供不确定性估计。

\section{未来工作}\label{sec:future}

基于上述局限性，本文提出以下未来工作方向。

\textbf{参数化统一PINN。}将当前固定参数的统一PINN扩展为真正的参数化版本，在参数空间$(\sigma,\beta,\kappa,\theta,\xi,\rho)$上进行联合训练，使单个网络能够对任意参数组合即时定价，无需重新训练。这需要更大的网络容量和更精细的训练策略（如课程学习）。

\textbf{A股市场适配。}当前框架的训练参数范围针对美股市场校准，A股市场（上证50ETF期权、沪深300ETF期权）的隐含波动率特征与美股有显著差异：A股波动率通常更高（年化波动率常超过30\%），且受政策因素影响，均值回归速度$\kappa$和长期均值$\theta$的分布与美股不同。未来可针对A股市场重新校准参数范围，并适配AkShare等国内数据接口，同时处理涨跌停限制、T+1交割等市场微结构特殊规则对定价的影响。

\textbf{美式期权扩展。}将统一框架扩展至美式期权，处理自由边界（最优执行边界）问题。美式期权的定价等价于最优停止问题，对应的PDE为变分不等式，需要引入惩罚函数法或互补条件。可以借鉴现有的PINN美式期权工作\cite{louskos2021american}，将自由边界条件纳入统一损失函数。

\textbf{奇异期权扩展。}障碍期权（Barrier Options）的定价PDE与欧式期权相同，但边界条件不同（在障碍处$V=0$或$V=\text{回扣}$），理论上可纳入本文框架，需额外设计障碍边界的网络处理方式。亚式期权（Asian Options）的收益依赖路径平均价格，引入路径依赖变量后维度更高，需要更大规模的网络和采样策略。

\textbf{Greeks计算。}利用PINN的自动微分能力，直接计算期权的Greeks（Delta、Gamma、Vega等），无需有限差分近似。这是PINN相对于传统数值方法的天然优势，在风险管理中具有重要价值。

\textbf{本地LLM部署。}将LLM路由层替换为本地部署的小型语言模型，降低API依赖和调用延迟，同时保护用户的期权参数隐私。

\textbf{实盘数据验证。}当前实验仅使用合成数据（参考解）进行验证，未来需要在真实市场期权数据上进行验证，评估统一PINN在实际市场条件下的定价精度。

\section{本章小结}\label{sec:conclusion_summary}

本文提出了基于PINN与LLM的统一期权定价框架，通过软掩码机制、加法参数化输出和LLM路由层，实现了BSM、CEV、Heston三大模型的统一求解和自然语言交互。实验结果表明，统一框架在精度和效率上均达到实用水平。本文的工作为PINN在金融工程中的应用提供了新的思路，也为LLM与科学计算工具的结合提供了一个具体的实践案例。
